% 本文件由 paperswarm.tex_gate 从台账自动生成，请勿手工编辑。
% 每个数值均由证据台账注入：claim_id -> (artifact SHA-256, JSON Pointer)。
\section{Method}

The method employed in PaperSwarm involves a multi-agent system that collaboratively generates a research paper through an evidence-gated process. The system evaluates two regression models: the ordinary least squares (OLS) and ridge regression. The OLS solution is obtained by minimizing the sum of squared residuals, while ridge regression introduces a penalty term to the loss function to prevent overfitting.

The mean test mean squared error (MSE) of the minimum-norm OLS solution, averaged over 8 seeds, is denoted by 0.6198. The standard deviation of the OLS test MSE across seeds is 0.1329, indicating high variance in the performance of OLS across different random initializations. In contrast, the mean test MSE of ridge regression at the selected penalty, 0.1, is 0.3041, with a lower standard deviation of 0.0623, suggesting more stable performance. The relative reduction in mean test MSE achieved by ridge over OLS is 50.94 percent, highlighting the effectiveness of ridge regression in reducing the irreducible noise floor, 0.1225, used as a reference.

The system is designed to handle datasets with 120 features and 90 training samples per split. Each configuration is evaluated over 8 random seeds to ensure robustness. The test samples per split are 200, providing a fair evaluation of the models' generalization capabilities.

The multi-agent system iteratively refines its outputs based on the evidence gathered from these experiments. The results indicate that ridge regression outperforms OLS by a significant margin, as evidenced by the lower mean test MSE and the reduced variance in performance. These findings support the use of ridge regression for improving the robustness and generalization of the models in PaperSwarm. The methodology is further validated by the consistent performance across multiple random initializations, as demonstrated by the low standard deviation of the ridge test MSE.
